\section{Hệ thống monolith cũ}
\label{ch:monolith}

\subsection{Kiến trúc và triển khai}

Monolith (nhánh \code{monolith}) là một ứng dụng ASP.NET Core (\code{Clinic.API}) chia lớp API / Application / Domain / Infrastructure, đi kèm frontend Vue và một SQL Server. Các nghiệp vụ gọi nhau trong cùng tiến trình và dùng một \code{ApplicationDbContext} duy nhất (Hình~\ref{fig:monolith}). Đây là khác biệt chính so với ranh giới sở hữu dữ liệu của bản microservices.

Để benchmark đối chứng, monolith được container hóa bằng \path{docker-compose.monolith.yml} (API ở cổng 5100) với tổng ngân sách CPU tương đương bản microservices (12,5 CPU).

\begin{figure}[!t]
\centering
\resizebox{\linewidth}{!}{%
\begin{tikzpicture}[node distance=9mm and 12mm]
  \node[actor] (u) {Patient / Receptionist\\Doctor / Admin};
  \node[box, right=of u] (fe) {Vue Frontend};
  \node[box, right=of fe, minimum width=4.2cm] (api) {\textbf{Clinic.API (monolith)}\\Controllers + Application\\+ Domain};
  \node[box, below=of api, minimum width=4.2cm] (inf) {Infrastructure / EF Core\\\code{ApplicationDbContext}};
  \node[db, below=of inf] (db) {SQL Server\\ClinicDB};
  \draw[arr] (u) -- (fe);
  \draw[arr] (fe) -- node[lbl,above]{HTTP} (api);
  \draw[arr] (api) -- node[lbl,right]{gọi trong tiến trình} (inf);
  \draw[arr] (inf) -- (db);
\end{tikzpicture}%
}
\caption{Kiến trúc hệ thống monolith}
\label{fig:monolith}
\end{figure}

\subsection{Thiết kế cơ sở dữ liệu}

\code{ApplicationDbContext} đăng ký chung các nhóm bảng:
\begin{codeitems}
  \item \textbf{Định danh và phân quyền:} \code{Persons}, \code{Users}, \code{Roles}, \code{Permissions}, \code{UserRoles}, \code{RolePermissions}, \code{RefreshTokens}.
  \item \textbf{Nhân sự và lịch làm việc:} \code{Employees}, \code{Doctors}, \code{Specialties}, \code{WorkHistories}, \code{WorkSchedules}, \code{ShiftRequests}.
  \item \textbf{Khám bệnh:} \code{Patients}, \code{Appointments}, \code{AppointmentSnapshots}, \code{QueueTickets}, \code{QueueCounters}.
  \item \textbf{Hồ sơ y tế:} \code{MedicalReports}, \code{MedicalReportAmendments}, \code{MedicalAccessAudits}.
  \item \textbf{Thông báo:} \code{Notifications}.
\end{codeitems}

Cấu hình EF Core nằm ở \path{backend/Clinic.Infrastructure.Sqlserver/Configurations}. Hình~\ref{fig:erd} là ERD khái quát các quan hệ nghiệp vụ chính (không liệt kê mọi cột, index hay migration).

\begin{figure*}[t]
\centering
\resizebox{\linewidth}{!}{%
\begin{tikzpicture}[
  ent/.style={draw=svcborder,fill=svc,rounded corners=1pt,font=\footnotesize\bfseries,
              minimum width=2.5cm,minimum height=7mm,align=center},
  rel/.style={draw=black!65,thick},
  card/.style={font=\tiny,fill=white,inner sep=1pt},
  node distance=11mm and 9mm]
  % Cột giữa: Person
  \node[ent] (person) {PERSON};
  \node[ent, above left=of person] (user) {USER};
  \node[ent, above=of person] (employee) {EMPLOYEE};
  \node[ent, below=of person] (patient) {PATIENT};
  \node[ent, left=of person] (notif) {NOTIFICATION};
  % Định danh
  \node[ent, above=of user] (rtoken) {REFRESH\_TOKEN};
  \node[ent, left=of user] (userrole) {USER\_ROLE};
  \node[ent, above=of userrole] (role) {ROLE};
  \node[ent, left=of role] (roleperm) {ROLE\_PERM};
  \node[ent, below=of roleperm] (perm) {PERMISSION};
  % Nhân sự
  \node[ent, right=22mm of employee] (doctor) {DOCTOR};
  \node[ent, above=of doctor] (workhist) {WORK\_HISTORY};
  \node[ent, right=of workhist] (specialty) {SPECIALTY};
  \node[ent, right=of doctor] (schedule) {WORK\_SCHEDULE};
  % Khám
  \node[ent, below=24mm of schedule] (appt) {APPOINTMENT};
  \node[ent, right=of appt] (snap) {APPT\_SNAPSHOT};
  \node[ent, below=of appt] (ticket) {QUEUE\_TICKET};
  \node[ent, below left=of ticket] (report) {MEDICAL\_REPORT};
  \node[ent, below=of report] (amend) {AMENDMENT};

  \draw[rel] (person) -- node[card]{1 : 0..1} (user);
  \draw[rel] (person) -- node[card]{1 : 0..1} (employee);
  \draw[rel] (person) -- node[card]{1 : 0..1} (patient);
  \draw[rel] (person) -- node[card]{1 : n} (notif);
  \draw[rel] (user) -- node[card]{1 : n} (rtoken);
  \draw[rel] (user) -- node[card]{1 : n} (userrole);
  \draw[rel] (role) -- node[card]{1 : n} (userrole);
  \draw[rel] (role) -- node[card]{1 : n} (roleperm);
  \draw[rel] (perm) -- node[card]{1 : n} (roleperm);
  \draw[rel] (employee) -- node[card]{1 : 0..1} (doctor);
  \draw[rel] (doctor) -- node[card]{1 : n} (workhist);
  \draw[rel] (specialty) -- node[card]{1 : n} (workhist);
  \draw[rel] (doctor) -- node[card]{1 : n} (schedule);
  \draw[rel] (schedule) -- node[card]{1 : n} (appt);
  \draw[rel] (patient) -- node[card]{1 : n} (appt);
  \draw[rel] (appt) -- node[card]{1 : n} (snap);
  \draw[rel] (appt) -- node[card]{1 : 0..1} (ticket);
  \draw[rel] (appt) -- node[card]{1 : 0..1} (report);
  \draw[rel] (ticket) -- node[card]{1 : 0..1} (report);
  \draw[rel] (report) -- node[card]{1 : n} (amend);
\end{tikzpicture}%
}
\caption{ERD khái quát của CSDL monolith}
\label{fig:erd}
\end{figure*}

Một số quan hệ là tùy chọn: Person có thể chưa có User; khách walk-in vẫn có hồ sơ Person/Patient. \code{MedicalAccessAudits} nằm cùng database nhưng chỉ khai báo khóa và index, không khai báo khóa ngoại trong EF nên không xuất hiện trên sơ đồ.

\subsection{Hạn chế của monolith}

\begin{itemize}
  \item Mọi miền nghiệp vụ chia sẻ một mô hình dữ liệu; thay đổi schema ở một miền có thể ảnh hưởng miền khác.
  \item Không thể mở rộng riêng phần chịu tải cao (đặt lịch, hàng đợi) mà phải nhân bản toàn bộ ứng dụng.
  \item Lỗi ở một module (ví dụ gửi thông báo) có thể làm chậm hoặc làm hỏng cả tiến trình.
  \item Phát hành phải triển khai lại toàn bộ ứng dụng.
\end{itemize}
